% Rutas y nombres en \texttt no se dividen: se deja holgura al espaciado.
\sloppy

\chapter{Register Map}
\label{app:regs}

This appendix is the interface contract between the engine and its host,
version 0x03 of the \ac{RTL}. The same text, with the history of each
rule, is kept in the repository as \nolinkurl{docs/mapa_registros.md}; the
\ac{RTL} (\nolinkurl{rtl/top/motor_top.vhd}), the host model
(\nolinkurl{firmware/host_check/motor_sim.c}) and the access layer of the
firmware (\nolinkurl{firmware/componentes/motor_i2c}) follow it.

The engine answers at 7-bit \ac{I2C} address 0x30. A write sets an 8-bit
register pointer that auto-increments on every data byte, read or
written. The slave never stretches the clock: long commands are launched
by writing \texttt{CONTROL} and the host polls \texttt{STATUS} until
\texttt{busy} clears. Undefined addresses read as 0x00 and writes to
read-only addresses are ignored. Register, bit and command names are
translated here; the repository uses the Spanish ones
(\texttt{CLAVE}, \texttt{RETO}, \texttt{sembrar}, \ldots).

\begin{table}[htbp]
\centering\small
\caption{Registers. R: read, W: write.}
\label{tab:app_regs}
\begin{tabular}{@{}llc>{\raggedright\arraybackslash}p{0.56\linewidth}@{}}
\toprule
Address & Name & Access & Content \\
\midrule
0x00 & ID & R & 0xA5 \\
0x01 & VERSION & R & 0x03 \\
0x02 & STATUS & R & Table~\ref{tab:app_status} \\
0x03 & CONTROL & W & commands, Table~\ref{tab:app_control}; self-clearing \\
0x04 & CFG\_KD & R/W & divider exponent, $K_D = 2^{\text{value}}$, effective
  range 0--24 (larger values act as 24). Resets to 24. Written only while
  the engine is idle; a write during a command is ignored \\
0x05 & CFG\_PAG & R/W & page of the capture buffer, 0--3 \\
0x06--0x07 & RCT\_MAX & R & longest run seen by the \ac{RCT}, low byte first \\
0x08--0x09 & APT\_MAX & R & largest count seen by the \ac{APT}, low byte first \\
0x0A--0x0D & REQUESTS & R & reseed counter of the \ac{DRBG}, 32 bits, low byte
  first. Saturates at $2^{32}-1$ \\
0x0E & DIAG & R & Table~\ref{tab:app_diag} \\
0x10--0x1F & KEY & R & current key, 16 bytes. \textbf{Test mode only} \\
0x20--0x2F & CHALLENGE & W & 16-byte nonce for challenge-response \\
0x30--0x3F & TAG & R & $\mathrm{CMAC}_{\text{KEY}}(\text{CHALLENGE})$, 16 bytes;
  reads as zeros while \texttt{tag\_ready} is 0 \\
0x40--0x5F & RANDOM & R & 32 bytes from the last \texttt{random} command, a
  \ac{DRBG} generation independent of the key \\
0x80--0xFF & CAPTURE & R & 128 bytes of raw bits, page CFG\_PAG.
  \textbf{Test mode only} \\
\bottomrule
\end{tabular}
\end{table}

\begin{table}[htbp]
\centering\small
\caption{STATUS (0x02).}
\label{tab:app_status}
\begin{tabular}{@{}cl>{\raggedright\arraybackslash}p{0.6\linewidth}@{}}
\toprule
Bit & Name & Meaning \\
\midrule
0 & seeded & the \ac{DRBG} holds a valid state \\
1 & busy & a command is running; do not send another \\
2 & rct\_alarm & repetition count test fired (stuck source) \\
3 & apt\_alarm & adaptive proportion test fired (biased source) \\
4 & key\_valid & a key is loaded in the authenticator \\
5 & capture\_full & the capture buffer holds a chunk \\
6 & tag\_ready & TAG belongs to the current CHALLENGE \\
7 & test\_mode & the test pin is active \\
\bottomrule
\end{tabular}
\end{table}

\begin{table}[htbp]
\centering\small
\caption{DIAG (0x0E).}
\label{tab:app_diag}
\begin{tabular}{@{}cl>{\raggedright\arraybackslash}p{0.6\linewidth}@{}}
\toprule
Bit & Name & Meaning \\
\midrule
0 & started & the start-up test has passed: one full \ac{APT} window (1024
  samples) without alarm since reset or since the last
  \texttt{clear\_alarms} \\
1 & rejected & the last command written to CONTROL was ignored, or, for a
  seeding, aborted. Cleared by the next accepted command \\
2 & exhausted & the request counter is at $2^{32}-1$; \texttt{generate} and
  \texttt{random} are refused until the engine is reseeded \\
\bottomrule
\end{tabular}
\end{table}

\begin{table}[htbp]
\centering\small
\caption{CONTROL (0x03). One command per write.}
\label{tab:app_control}
\begin{tabular}{@{}cl>{\raggedright\arraybackslash}p{0.66\linewidth}@{}}
\toprule
Bit & Command & Action \\
\midrule
0 & seed & collect 384 raw bits (256 of entropy, 128 of nonce), wait until
  the health tests have judged all of them, and instantiate the \ac{DRBG}
  from scratch \\
1 & generate & draw 256 bits; the first 128 become KEY and are loaded into
  the authenticator, the rest are discarded. Nothing becomes readable \\
2 & authenticate & TAG $\leftarrow \mathrm{CMAC}_{\text{KEY}}(\text{CHALLENGE})$ \\
3 & clear\_alarms & clear the health alarms; restarts the start-up window \\
4 & capture & fill the capture buffer with raw bits. Test mode only, engine
  idle, no capture in progress \\
5 & reseed & as \texttt{seed}, but on top of the current state; requires a
  seeded \ac{DRBG} \\
6 & random & draw 256 bits into RANDOM \\
7 & clear\_all & zeroization: KEY, the whole \ac{DRBG} state and its
  intermediates (including the \ac{AES} core), CHALLENGE, RANDOM, TAG and
  the collected seed to zero, and the zero key loaded into the
  authenticator. Accepted whenever the engine is idle, also with an
  alarm \\
\bottomrule
\end{tabular}
\end{table}

\paragraph{Rules enforced by the hardware.} A refused command does not
raise \texttt{busy} and sets \texttt{DIAG.rejected}.
\begin{itemize}
\item With a health alarm set, \texttt{seed}, \texttt{reseed},
  \texttt{generate} and \texttt{random} are refused; \texttt{authenticate}
  is still served, because it uses the key already generated and draws
  no entropy.
\item No seeding before the start-up test has passed
  (SP~800-90B~\S4.3~\cite{turan_2018_sp80090b}). The seed enters the
  \ac{DRBG} only after the \ac{APT} window that contains its last bit
  closes without alarm; an alarm or a \texttt{clear\_alarms} during
  collection aborts the seeding and leaves the previous state untouched.
  A seeding takes at most $384 + 1024$ raw bits.
\item \texttt{seed} and \texttt{reseed} are refused with
  \texttt{CFG\_KD} below the minimum (generic \texttt{G\_KD\_MIN}, 6 by
  default), below which the bit-level clock-domain crossing loses bits
  silently and the health tests cannot tell.
\item A capture and a seeding exclude each other, so the capture buffer
  never holds bits of a seed.
\item The request counter does not wrap: at $2^{32}-1$ the engine
  refuses to generate until reseeded (\texttt{DIAG.exhausted}). The
  reseeding policy below that ceiling is the host's.
\end{itemize}

\paragraph{Priority.} If several command bits arrive in one write, only
the highest-priority one is served: \texttt{clear\_all}, then
\texttt{seed}/\texttt{reseed} (\texttt{reseed} if both), \texttt{generate},
\texttt{random}, \texttt{authenticate}, \texttt{capture}.
\texttt{clear\_alarms} is served alongside any of them, but it must not
share a write with a seeding: it restarts the start-up window, so the
seeding is accepted and aborted on the next cycle.

\paragraph{Typical host sequence.}
\begin{enumerate}
\item Read ID (0xA5) and VERSION (0x03).
\item Write \texttt{CFG\_KD} with the divider of the characterization, or
  leave the reset value.
\item Poll DIAG until \texttt{started}; write \texttt{seed}; poll STATUS
  until \texttt{busy} clears; check \texttt{DIAG.rejected} $= 0$,
  \texttt{seeded} $= 1$ and no alarm.
\item Write \texttt{generate}, poll, check \texttt{DIAG.rejected}. A key
  now exists.
\item Bench provisioning only, in test mode: read KEY.
\item Challenge-response: write CHALLENGE, write \texttt{authenticate},
  poll, read TAG and compare it with the \ac{CMAC} computed by the host.
\item Random for the host: write \texttt{random}, poll, read RANDOM.
\end{enumerate}

\paragraph{Test mode.} The \texttt{test\_mode} pin, synchronized by two
flip-flops, enables reading KEY and CAPTURE and the \texttt{capture}
command. On the \ac{FPGA} it is a board switch; on an integrated circuit
it would be a fuse blown at production.

\chapter{Official Test Vectors Used}
\label{app:vectors}

Table~\ref{tab:app_vectors} lists where every reference value used by the
benches comes from, and how each check closes. Official vectors are used
where a standard publishes them. Where the hardware's input layout
differs from any published case, the expected values come from a
reference model in Python that is itself checked against the official
vectors first, and written independently of the \ac{VHDL}.

\begin{table}[htbp]
\centering\small
\caption{Reference values and their origin. Counts are checks passed in
the last run (\nolinkurl{results/tb_motor_top.log}, 177 in total, with
11 of 11 \acs{PSL} assertions and 6 of 6 cover points).}
\label{tab:app_vectors}
\begin{tabular}{@{}>{\raggedright\arraybackslash}p{0.2\linewidth}>{\raggedright\arraybackslash}p{0.4\linewidth}>{\raggedright\arraybackslash}p{0.32\linewidth}@{}}
\toprule
Block & Source of the expected values & Where checked \\
\midrule
\ac{AES}-128 (both cores) & FIPS~197~\cite{nist_2023_fips197},
  Appendix~B and Appendix~C.1; SP~800-38A, F.1.1 ECB-AES128, three
  blocks; a thousand chained encryptions repeated; Appendix~C.1 again
  after \texttt{clear} & \texttt{tb\_aes\_enc}, 18 checks, the two cores
  with the same stimuli and required to agree \\
\ac{AES}-\ac{CMAC} & RFC~4493~\cite{song_2006_rfc4493} (the examples of
  SP~800-38B~\cite{nist_2016_sp80038b}): the four messages of 0, 16, 40
  and 64 bytes under key \texttt{2b7e1516\ldots}, and the 16-byte case
  repeated; the subkeys $K_1$ and $K_2$ of the same example & \texttt{tb\_aes\_cmac}, 7 checks \\
\ac{CTRDRBG}, reference model &
  \ac{CAVP} \texttt{drbgtestvectors}~\cite{nist_cavp_drbg},
  \texttt{CTR\_DRBG\_no\_reseed.rsp} and \texttt{CTR\_DRBG\_pr\_false.rsp},
  AES-128 with and without derivation function: $4 \times 240 = 960$ cases &
  \nolinkurl{analysis/ctr_drbg_ref.py}, 960 of 960 \\
\ac{CTRDRBG}, \ac{RTL} & the reference model above, with the hardware's
  input layout (256-bit entropy, 128-bit nonce, no personalization):
  \nolinkurl{analysis/gen_drbg_vectores.py} writes
  \texttt{drbg\_vectores\_pkg.vhd} & \texttt{tb\_ctr\_drbg}, 18 checks,
  including zeroization and counter saturation \\
Whole engine & the key comes from ring noise, so the bench dumps key,
  challenge and tag, and the \ac{CMAC} is recomputed with the Python model
  & \nolinkurl{analysis/verifica_top.py}, 5 of 5 \\
Firmware against the engine model & \ac{CMAC} and ECDSA of the target's
  mbedtls 3.6, compiled for the PC; signed readings checked with the
  \texttt{cryptography} library & \texttt{run\_host.py}, 0 failures;
  \texttt{verifica\_lecturas.py} \\
Health-test cut-offs & SP~800-90B~\S4.4~\cite{turan_2018_sp80090b}
  formulas for $H = 0.98$ and $\alpha = 2^{-20}$: \ac{RCT} $C = 22$, \ac{APT} $W = 1024$,
  $C = 596$ & \nolinkurl{analysis/cortes_salud.py}; behaviour in
  \texttt{tb\_health\_tests}, 15 checks \\
\bottomrule
\end{tabular}
\end{table}

The CAVP files are kept in the repository
(\nolinkurl{tests/vectors/drbg/}), so the validation does not depend on the
\ac{NIST} site staying as it is. The entropy estimates of
Chapter~\ref{ch:ch7_results} are not test vectors: they come from the
\ac{NIST} SP~800-90B tool~\cite{nist_ea_github}, run unchanged.

\chapter{Reproduction Guide}
\label{app:repro}

Everything in this work that does not need the boards is rerun on every
push by a continuous-integration workflow
(\nolinkurl{.github/workflows/verifica.yml}) on a clean Ubuntu machine that
is not the author's: seven jobs, no credentials. The repository is
public: \url{https://github.com/mmggjj/TFM_ENCRY}. The commands below are
the ones the workflow runs, so a green run reproduces them. Outputs are
never filtered: the scripts print the full tool output and fail on any
error.

\begin{table}[htbp]
\centering\small
\caption{What the workflow reproduces.}
\label{tab:app_ci}
\begin{tabular}{@{}>{\raggedright\arraybackslash}p{0.27\linewidth}>{\raggedright\arraybackslash}p{0.38\linewidth}>{\raggedright\arraybackslash}p{0.27\linewidth}@{}}
\toprule
Job & Command & Result \\
\midrule
Model and estimators & \texttt{python} \nolinkurl{analysis/ero_model.py},
  \texttt{jitter\_estimator.py}, \texttt{ctr\_drbg\_ref.py},
  \texttt{cortes\_salud.py}, \texttt{fuzz\_verifica\_lecturas.py},
  \texttt{verifica\_lecturas.py} on the board log & self-checks, 960
  \ac{CAVP} cases, board readings valid \\
\ac{VHDL} benches & \texttt{python} \nolinkurl{rtl/run_tb.py}, then
  \texttt{python} \nolinkurl{analysis/verifica_top.py} & 7 benches, \acs{PSL}
  report as artefact \\
Synthesis & \texttt{python} \nolinkurl{synth/sintetiza.py} \texttt{--verilog} &
  area per block in kGE (Chapter~\ref{ch:ch6_asic}) \\
Firmware on the PC & \texttt{python} \nolinkurl{firmware/host_check/run_host.py}
  \texttt{--mbedtls <dir>}, and again with \texttt{--sanear} & 0 failures, with
  ASan and UBSan \\
ESP-IDF build & \texttt{idf.py set-target}, \texttt{build}, \texttt{size},
  and \texttt{idf.py -B build\_sim -DMOTOR\_SIMULADO=1 build}, for the
  verifier (ESP32) and the meter (ESP32-C3) & 4 images \\
Thesis & \texttt{tectonic tfm\_main.tex} and \texttt{tfm\_main\_es.tex} &
  both PDFs as artefacts \\
\bottomrule
\end{tabular}
\end{table}

\paragraph{Tools.} Python~3.12 with \texttt{numpy}, \texttt{scipy} and
\texttt{cryptography}; GHDL from the Ubuntu packages (5.0.1 with the mcode back-end on the
author's machine) for the benches, which must be run with \acs{VHDL}-2008; Yosys through
\texttt{yowasp-yosys}, or the OSS CAD Suite with the GHDL plug-in, and the
IHP SG13G2 standard-cell library of the open
\ac{PDK}~\cite{ihp_openpdk_repo}; mbedtls~3.6.2 from its release tarball
(a git clone lacks the generated sources); ESP-IDF~5.3.5, the workflow
using the \nolinkurl{espressif/idf:v5.3.5} container; tectonic~0.15. The
GHDL~7.0 development build with the LLVM back-end crashes on the
\acs{PSL} properties of \texttt{tb\_motor\_top} that use external names,
and should not be used.

\paragraph{Outside the workflow.} Three parts need software or hardware
the workflow does not have:
\begin{itemize}
\item \emph{Transistor-level ring.} ngspice~47~\cite{ngspice_47} with
  OpenVAF-Reloaded~\cite{openvaf_reloaded} to compile the PSP models of
  the \ac{PDK}; \texttt{python} \nolinkurl{analysis/anillo_ihp_repeticiones.py},
  with \texttt{--esquina}, \texttt{--temp}, \texttt{--vdd} and
  \texttt{--mismatch} for corners, supply and mismatch. ngspice cannot
  fix the noise realization, so every run is a new one and the
  uncertainty comes from repeating. One 400~ns realization takes about
  30 minutes.
\item \emph{SP~800-90B estimates of the ESP32.} The \ac{NIST}
  tool~\cite{nist_ea_github} built from source;
  \texttt{python} \nolinkurl{analysis/entropia_90b.py}, then
  \texttt{resumen\_90b.py}, on the captures kept in \nolinkurl{results/};
  the restart test with \texttt{reinicios\_90b.py}. New captures need an
  ESP32 and \nolinkurl{analysis/campana_rng.py}.
\item \emph{\ac{FPGA}.} Vivado (free tier, Artix-7 included) and
  \nolinkurl{synth/basys3/build.tcl}; installation and programming are
  described in \nolinkurl{docs/instalacion_vivado_windows.md}.
\end{itemize}

\paragraph{First ESP32 campaign.} The analysis of the first campaign
(Sections~\ref{sec:esp32_sp80022} to~\ref{sec:esp32_harvester}) reruns
without the board from the files kept in
\nolinkurl{results/esp32_rng_agosto}, with \texttt{numpy}, \texttt{scipy},
\texttt{matplotlib} and \texttt{cryptography}:
\begin{itemize}
\item \texttt{python} \nolinkurl{analysis/esp32_rng/nist_sp800_22.py} runs
  the self-test of the SP~800-22 battery (operating-system data, zeros,
  alternating pattern).
\item \texttt{python} \nolinkurl{analysis/esp32_rng/harvester_from_sources.py}
  writes the 1000 seeds of the off-line validation to
  \texttt{seeds\_rep1.bin}, from \texttt{puf\_boot1.bin} and
  \texttt{condD\_rep1.bin}.
\item \texttt{python} \nolinkurl{analysis/esp32_rng/run_campaign.py} runs
  \texttt{nist\_analysis.py}, \texttt{entropy\_analysis.py},
  \texttt{temperature\_correlation.py}, \texttt{harvester\_validation.py}
  and the off-line mode of \texttt{mqtt\_tls\_demo.py}. Each writes its CSV
  and figures into the data directory, and \texttt{run\_campaign.py} copies
  the figures into \nolinkurl{memoria/figures} under their own names; the
  thesis uses only \texttt{nist\_heatmap.pdf}, as
  \texttt{esp32\_nist\_heatmap.pdf}. The four scripts also run one by one,
  with \texttt{-d <data directory>} and \texttt{--figdir <directory>}.
  Rerun on 2026-10-31 on a copy of the data, they reproduce the CSV files
  kept, value for value.
\item \texttt{python} \nolinkurl{analysis/esp32_rng/make_result_tables.py}
  writes the result tables as \LaTeX{} into \texttt{tablas\_generadas.tex},
  and \texttt{python} \nolinkurl{analysis/esp32_rng/figuras_memoria.py}
  draws Figures~\ref{fig:esp32_metrics}, \ref{fig:esp32_temp}
  and~\ref{fig:esp32_harvester} into \nolinkurl{memoria/figures}.
\end{itemize}
A new capture needs an ESP32 with \nolinkurl{firmware/esp32_rng} built and
flashed (\texttt{idf.py set-target esp32}, \texttt{build},
\texttt{-p <port> flash}; ESP-IDF~5.3, console at 115\,200~baud). Then
\texttt{python} \nolinkurl{analysis/esp32_rng/capture_rng.py}
\texttt{--port <port>} \texttt{--reps <n>} resets the board through RTS
and collects $n$ valid sequences per condition, with the SRAM patterns
(\texttt{puf\_boot<k>.bin}), the seeds and \texttt{capture\_log.csv}; by
default it writes into \nolinkurl{results/esp32_rng_agosto}, so
\texttt{-o <directory>} keeps new data apart. A reset through RTS keeps
the SRAM powered and does not give a new power-up pattern: for the clean
series, unplug the board, plug it in again and run
\texttt{python} \nolinkurl{analysis/esp32_rng/capture_puf.py}
\texttt{--port <port>} once per power-up, which adds one
\texttt{puf\_pc<N>.bin}. The later campaigns use the same firmware with
\nolinkurl{analysis/campana_rng.py}, which keeps the whole serial stream
and numbers the repetitions from \texttt{--primera-rep}.
